\documentclass[a4paper,12pt]{article}
\usepackage[utf8]{inputenc}
\usepackage[T1]{fontenc}
\usepackage[norsk]{babel}
\usepackage{amsmath}
\usepackage{listings}
\usepackage{xcolor}
\usepackage{geometry}
\geometry{margin=2.5cm}
\usepackage{enumitem}
\usepackage[T1]{fontenc}
\usepackage{lmodern}
\usepackage{cmap}
\usepackage{needspace}
\usepackage[colorlinks=true, linkcolor=red!50!black, citecolor=red!50!black, urlcolor=blue]{hyperref}

\usepackage{setspace}
\onehalfspacing

\lstset{
    language=Python,
    basicstyle=\ttfamily\small,
    keywordstyle=\color{blue},
    commentstyle=\color{green!60!black},
    stringstyle=\color{red},
    showstringspaces=false,
	columns=fixed,
    keepspaces=true,
    upquote=true,
	inputencoding=utf8,
    frame=single,
	tabsize=4,
    breaklines=false,
	    literate=
     {æ}{{\ae}}1
     {ø}{{\o}}1
     {å}{{\aa}}1
     {Æ}{{\AE}}1
     {Ø}{{\O}}1
     {Å}{{\AA}}1
}

\title{Eulers metode i Python}
\author{}
\date{}

\begin{document}

\maketitle

\section*{Mål}
Målet med oppgaven er å forstå og programmere \textit{Eulers metode} for å finne en numerisk tilnærming til løsningen av en differensialligning.

\section*{Teori}

En differensialligning på formen

\begin{subequations}
\begin{align}
&\frac{dy}{dx} = f(x,y) \\
\implies &dy = f(x,y) \, dx
\label{eq:dy}
\end{align}
\end{subequations}
beskriver hvordan den deriverte til en funksjon avhenger av både $x$ og $y$. Den kan være så enkel som
\[
f(x,y) = x,
\]
eller så komplisert som
\begin{equation}
\label{eq:vanskelig}
f(x,y) = \arctan{\left( \frac{e^{xy}}{x^2+y^2} \right)}.
\end{equation}
Vanskelige utrykk er like enkle å løse numerisk som enkle utrykk!
\newline

Dersom vi kjenner startpunktet (initsialbetingelsen) $(x_0,y_0)$ kan vi bruke \textit{Eulers metode} til å tilnærme løsningen numerisk. Neste punkt $(x_1,y_1)$ er gitt ved
\begin{subequations}
\label{eq:step1}
\begin{align}
x_1 &= x_0 + dx \\
\intertext{og}
y_1 &= y_0 + dy \\
\implies y_1 &= y_0 + f(x_0,y_0)\, dx \quad (\text{fra \eqref{eq:dy}}).
\end{align}
\end{subequations}
Mønstret fortsetter slik:
\begin{subequations}
\begin{align}
x_{2\phantom{+1}}     &= x_1 + dx, \\
y_{2\phantom{+1}}     &= y_1 + f(x_1,y_1)\, dx,\\
&\, \, \, \vdots  \notag\\
x_{n+1} &= x_n + dx, \\
y_{n+1} &= y_n + f(x_n,y_n)\, dx.
\end{align}
\end{subequations}

\section*{Oppgave 1: Forstå metoden}

Vi ser på differensialligningen
\begin{equation}
\frac{dy}{dx} = x + y,
\label{eq:diff1}
\end{equation}
eller med andre ord har vi \[f(x,y) = x + y.\]
Initialbetingelsen er gitt ved
\[
y(0) = 1,
\]
eller med andre ord er startpunktet $(x_0,y_0) = (0,1)$.
\newline

Du skal nå regne ut de første stegene manuelt for dette tilfellet. Bruk $dx=0.1$.

\begin{enumerate}[label=\alph*)]
\item Regn ut $x_1$ og $y_1$ ved hjelp av formelen \eqref{eq:step1}.\footnote{$(x_1,y_1) = (0.1, 1.1)$.}
\item Regn ut $(x_2,y_2)$.\footnote{$(x_2,y_2) = (0.2, 1.22)$.}
\item Fortsett til du går lei og konkluderer med at dette lar seg enklere gjøre i Python.\footnote{$(x_3,y_3) = (0.3, 1.362)$.\newline$\phantom{mmmmmiii} \vdots$}
\end{enumerate}
\newpage

\section*{Oppgave 2: Parametere}

Start med å definerer parameterene:

\begin{lstlisting}
x0 = 0
x_slutt = 1
y0 = 1
dx = 0.1
\end{lstlisting}

\section*{Oppgave 3: Programmer funksjonen}

Definer funksjonen $f(x,y)$ i Python:

\begin{lstlisting}
def f(x, y):
    return x + y
\end{lstlisting}


\section*{Oppgave 4: Implementer Eulers metode}

Skriv programmet:

\begin{lstlisting}
x = x0
y = y0

while x <= x_slutt:
    y = y + f(x, y) * dx
    x = x + dx
    print(x, y)
\end{lstlisting}

Svar på spørsmålene:

\begin{enumerate}[label=\alph*)]
\item Hvilken linje i programmet tilsvarer formelen 
\[
y_{n+1} = y_n + f(x_n,y_n)\, dx?
\]
\item Kjør programmet og sjekk at utregningene fra oppgave 1 stemmer.
\item Gjør Python noe dumt?
\end{enumerate}

\section*{Oppgave 5: Tegn grafen}

Utvid programmet slik at du lagrer $x$- og $y$-verdiene i lister og tegner grafen ved hjelp av \lstinline|matplotlib|.

\begin{lstlisting}
x_verdier = [x0]
y_verdier = [y0]

x = x0
y = y0

while x <= x_slutt:
    y = y + f(x, y) * dx
    x = x + dx
	
    x_verdier.append(x)
    y_verdier.append(y)

import matplotlib.pyplot as plt

plt.plot(x_verdier, y_verdier, "-o")
plt.xlabel("x")
plt.ylabel("y")
plt.title("Tilnærming med Eulers metode")
plt.grid()
plt.show()
\end{lstlisting}

\section*{Oppgave 6: Sammenlign med eksakt løsning}

Differensialligningen har eksakt løsning
\begin{equation}
y = 2e^{x} - x - 1
\label{eq:yexact}
\end{equation}

\begin{enumerate}[label=\alph*)]
\item Sjekk at dette stemmer ved å derivere $y$ og vis at 
\[
\frac{dy}{dx} = y + x = 2e^{x} - 1
\]

\item Definer funksjonen $y(x)$ i programmet \eqref{eq:yexact}:
\begin{lstlisting}
import numpy as np

def y(x):
    return 2*np.exp(x) - x - 1
\end{lstlisting}

\item Legg til følgende kode for å plotte $y(x)$:
\Needspace{3\baselineskip}
\begin{lstlisting}
X = np.linspace(x0, x_slutt, 100)

plt.plot(X, y(X), "--")
\end{lstlisting}

\item Undersøk hva som skjer når du velger mindre verdier for $dx$. Det er lurt å bytte ut startverdiene i toppen av koden med følgende:
\begin{lstlisting}
x0 = 0
x_slutt = 2
y0 = 1
dx = 0.1
\end{lstlisting}
og endre:
\begin{lstlisting}
X = np.linspace(x0, x_slutt, 100)
\end{lstlisting}

\item Hva ville skjedd om du endret $dx$ uten å velge den nye definisjonen av \lstinline|antall_steg|?

\end{enumerate}

\section*{Bonus: Løs en uløselig differensialligning}

Du skal nå løse denne differensialligningen \eqref{eq:vanskelig}
\begin{equation*}
\frac{dy}{dx} = f(x,y) = \arctan{\left( \frac{e^{xy}}{x^2+y^2} \right)}.
\end{equation*}
Alt du må gjøre er å bytte ut \lstinline|f(x, y)| med:
\begin{lstlisting}
def f(x, y):
    return np.arctan(np.exp(x*y)/(x**2+y**2))
\end{lstlisting}
Det kan være lurt å endre på noen parametere:
\begin{lstlisting}
x0 = 0
x_slutt = 5
y0 = 0.1
dx = 0.0001
\end{lstlisting}
Og også kommentere ut den eksakte løsningen fra forrige oppgave siden den ikke er løsningen for denne oppgaven:
\begin{lstlisting}
#plt.plot(X, y(X), "--")
\end{lstlisting}

\end{document}
